\documentclass[10pt,twocolumn]{article}
\usepackage[T1]{fontenc}
\usepackage[utf8]{inputenc}
\usepackage{lmodern}
\usepackage[margin=1.8cm,top=2.4cm,bottom=2cm,columnsep=0.6cm]{geometry}
\usepackage{fancyhdr}
\usepackage{titlesec}
\usepackage{booktabs}
\usepackage{amsmath,amssymb,amsfonts}
\usepackage{microtype}
\usepackage{xcolor}
\usepackage{mdframed}
\usepackage{parskip}
\usepackage{enumitem}
\usepackage{hyperref}

\definecolor{rulered}{RGB}{180,30,30}
\definecolor{boxgray}{RGB}{245,245,245}
\definecolor{keywordgray}{RGB}{100,100,100}

\pagestyle{fancy}
\fancyhf{}
\fancyhead[L]{\textsc{\small Claude Mag}}
\fancyhead[C]{\small\textit{Machine Learning Research Digest}}
\fancyhead[R]{\small 2026-09-16}
\fancyfoot[C]{\small\thepage}
\renewcommand{\headrulewidth}{0.8pt}

\titleformat{\section}{\normalsize\bfseries\color{rulered}}{}{0pt}{}%
  [\vspace{-4pt}\color{rulered}\rule{\columnwidth}{0.4pt}]
\titlespacing{\section}{0pt}{8pt}{4pt}

\hypersetup{colorlinks=true,urlcolor=rulered,linkcolor=black}

\begin{document}
\setlength{\parindent}{0pt}
\setlength{\parskip}{4pt}

{\small\color{keywordgray}\textbf{Keywords:}
  tabular foundation models \textbullet{}
  FP8 quantization \textbullet{}
  in-context learning \textbullet{}
  attention speedup \textbullet{}
  TabPFN}\par\vspace{4pt}

{\LARGE\bfseries\raggedright
  Attention Quantization for Tabular
  Foundation Models}\par\vspace{4pt}

{\large\itshape\raggedright
  Tabular FM KV Matrices Scale Quadratically With Training Rows ---
  FP8 Quantization With Test-Training Row Alignment Achieves
  1.7$\times$ Speedup at Zero Accuracy Cost}\par\vspace{6pt}

{\small\textbf{By} Jonas M.\ K\"ubler, Benjamin J\"ager, Klemens Fl\"oge,
  Noah Hollmann, Frank Hutter
  (University of Freiburg; ELLIS Institute T\"ubingen)}\par\vspace{2pt}

{\small\color{keywordgray}arXiv:2609.13031}
\par\vspace{2pt}\noindent{\color{rulered}\rule{\columnwidth}{0.6pt}}\vspace{6pt}

Tabular foundation models such as TabPFN and TabICL perform in-context
learning on new datasets by encoding the entire training set as a
sequence of tokens and attending over all rows in a single forward pass.
Their attention matrices scale quadratically with training set size ---
a table with 10,000 rows and 4 attention heads contains 400 million
attention entries per layer, dwarfing the KV caches of generative LLMs.
FP8 quantization has transformed LLM serving, but applying it to tabular
FMs requires solving a unique alignment problem: the quantization error
on test rows must match that on training rows, or prediction accuracy
collapses.

\begin{mdframed}[backgroundcolor=boxgray,linecolor=rulered,linewidth=1pt,
    innerleftmargin=6pt,innerrightmargin=6pt,
    innertopmargin=6pt,innerbottommargin=6pt]
\textbf{\small AT A GLANCE}\par\vspace{4pt}
\begin{description}[leftmargin=0pt,labelsep=4pt,itemsep=2pt,parsep=0pt,
    font=\small\bfseries]
  \item[Problem:] Tabular FM KV matrices scale $O(n^2)$ with training
    rows; FP8 quantization without row-group alignment degrades
    accuracy by 3--5 points.
  \item[Why it matters:] At $n=10{,}000$ training rows, FP16 attention
    costs 800\,MB per layer; FP8 halves this and doubles bandwidth.
  \item[Method:] Dynamic per-group FP8 scale calibration with test/train
    row alignment; custom asymmetric-tile Triton kernel.
  \item[Result:] 1.7$\times$ speedup on H100 with $<$0.2-point
    accuracy change on TabPFN-v3 and TabICLv2.
  \item[Evidence:] TabArena (47 datasets), BeyondArena (21 datasets);
    ablations on alignment, format, tile shape.
\end{description}
\end{mdframed}

\section{The Big Picture}

In a tabular FM, all $n$ training rows and $m$ test rows are encoded as
tokens and processed in one forward pass. The attention matrix $A \in
\mathbb{R}^{(n+m)\times(n+m)}$ has $(n+m)^2$ entries per head per
layer. For $n = 10{,}000$, $m = 1{,}000$, and 4 heads at FP16:
\[
  \text{memory} = 4\times(1.1\times10^4)^2\times 2\,\text{B}
  \approx 968\,\text{MB per layer}
\]
This exceeds the KV cache of a 7B LLM processing a 128k-token context
at the same precision. FP8 quantization halves this cost and doubles
effective bandwidth utilization.

\section{The Quantization Error Misalignment Problem}

Standard FP8 quantization assigns a single scale $s$ to a tensor $X$:
\[
  s = \frac{\|X\|_\infty}{448}, \quad
  \hat{X} = \mathrm{round}(X/s)\cdot s
\]
(448 is the maximum value in FP8 e4m3fn format.)

In a tabular FM, the query matrix $Q = [Q_\text{train};\,Q_\text{test}]$
has two row groups with different distributions: training rows include
label information in their embeddings; test rows do not.
If $\|Q_\text{train}\|_\infty \gg \|Q_\text{test}\|_\infty$, then $s$
is calibrated to training rows, and test queries use only the low end
of the FP8 range --- far fewer distinct values --- introducing systematic
bias in test attention scores.

\section{Core Mechanism: Per-Group Dynamic Alignment}

SAS applies separate dynamic scales per row group:
\[
  s_g = \frac{\|Q_g\|_\infty}{448}, \quad g\in\{\text{train},\text{test}\}
\]
Computed from the tensor itself at each forward pass (no calibration
dataset). This equalizes relative quantization error:
\[
  \frac{\epsilon_q}{\|Q_g\|_\infty} \leq \frac{1}{2\times 448}
  \approx 0.001
\]
for both row groups, eliminating the systematic bias.

\section{Technical Formulation}

\textbf{FP8 e4m3fn:} 4 exponent bits, 3 mantissa bits, 1 sign bit.
Representable range $[-448, 448]$; step size $s/8$ per scale.

\textbf{Per-group quantized attention:}
\[
  \hat{Q}_g = \mathrm{round}(Q_g/s_g)\cdot s_g
\]
\[
  A_{ij} = \mathrm{softmax}\!\left(\frac{\hat{\mathbf{q}}_i
    \hat{\mathbf{k}}_j^\top}{\sqrt{d}}\right),\quad
  \mathbf{O}_i = \sum_j A_{ij}\hat{\mathbf{v}}_j
\]

\textbf{Alignment overhead:} One pass over the input to compute
per-group max values; $<$1\% of total attention compute time.

\textbf{Speedup bound:} Bandwidth-bound regime (large $n$):
\[
  \text{Speedup}(n) \approx \frac{16}{8}\cdot\frac{B_\text{FP16}}{B_\text{FP8}}
  = 2\times\left(1 - \epsilon_\text{kernel}\right)
\]
with kernel overhead $\epsilon_\text{kernel} \approx 15\%$ for the
custom Triton kernel.

\section{Experimental Results}

Tested on TabPFN-v3 and TabICLv2 on an H100 GPU:

\begin{center}
\begin{tabular}{llcc}
\toprule
Model & Precision & Speedup & Acc.\ Change \\
\midrule
TabPFN-v3 & FP16            & 1.00$\times$ & --- \\
TabPFN-v3 & FP8 (no align)  & 1.65$\times$ & $-$3.1\% \\
TabPFN-v3 & \textbf{FP8 (aligned)} & \textbf{1.70}$\times$ & $-$0.1\% \\
\midrule
TabICLv2  & FP16            & 1.00$\times$ & --- \\
TabICLv2  & FP8 (no align)  & 1.62$\times$ & $-$4.7\% \\
TabICLv2  & \textbf{FP8 (aligned)} & \textbf{1.68}$\times$ & $-$0.2\% \\
\bottomrule
\end{tabular}
\end{center}

Alignment eliminates 3--5 point accuracy degradation; the residual
$<$0.2 points is within noise. Benchmarks: TabArena (47 classification
datasets), BeyondArena (21 regression datasets).

\section{Ablations and Interpretation}

\textbf{Without per-group scale (global scale):} Accuracy drops
3.1--4.7 points --- the distribution mismatch between test and training
rows is the dominant failure mode.

\textbf{E4m3fn vs.\ E5m2:} E4m3fn (more mantissa bits) outperforms
E5m2 (more exponent bits) because attention logit distributions are
well-concentrated and benefit more from mantissa precision.

\textbf{Tile shape:} Custom asymmetric tiling adds 18\% speedup over
reusing a standard square-tile FP8 kernel, justified by the
asymmetric $m\times n$ ($m \ll n$) shape of tabular FM attention.

\textbf{Scaling with $n$:} Speedup grows from 1.3$\times$ at
$n=500$ to 1.9$\times$ at $n=10{,}000$, reflecting that larger
attention matrices are more bandwidth-bound.

\section*{Reference}

Jonas M.\ K\"ubler, Benjamin J\"ager, Klemens Fl\"oge, Noah Hollmann,
Frank Hutter. \textbf{Attention Quantization for Tabular Foundation
Models.} arXiv:2609.13031, September 2026.
\url{https://arxiv.org/abs/2609.13031}

\end{document}
